\section{An actual source-curvature payment when the append window crosses the old peak}
This is the peak-crossing sublemma subsequently used in the COMPLETE
triangular source-curvature theorem in
\texttt{triangular\_source\_curvature\_complete.tex}.
This sublemma alone does not prove general C, S, A2, or original UI.

Let $F$ be an actual nonrectangular sorted arm, with first missing weight
$\kappa$, leading source multiplicity $\delta$, and
\[
 E=1-z+(2z-1)F,\qquad D_n=(F_n-F_{n-1})^+\quad(n\ge2).
\]
Put $D_0=D_1=0$, let $\ell$ be the FIRST global mode of $D$, and
$M=D_\ell$. Append a legal capacity $t$, and use hats for the new arm.
Its leading source multiplicity is $\widehat\delta\le\widehat\kappa
\le\kappa+1$. All sequences are zero extended.

\begin{theorem}[Peak-crossing source payment]
If $n$ is at or before the FIRST global mode of $\widehat D$ and
$n-1>\ell$, then
\[
                    \Delta\widehat E_n\ge M-1.                 \tag{1}
\]
If moreover $\kappa\ge8$, then
\[
 \Delta\widehat E_n\ge
       \max\{1,\tbinom{\widehat\delta+1}{2}-1\}.                \tag{2}
\]
In particular, the peak-crossing portion of the proposed triangular
source-curvature bound is unconditionally paid for all old
$\kappa\ge8$, with no bound on initial ones, capacities, or count.
\end{theorem}

\begin{proof}
Let $m$ be the LAST ordinary mode of $F$. The established positive
derivative splice says that the new derivative agrees, through rank
$m+1$, with
\[
                     P=z^2[t-1]+S_tD,
\]
and is nonincreasing thereafter. This also gives the joining edge;
therefore the first new derivative mode is at most $m+1$.

Since $D$ is nonincreasing from its first mode $\ell$, the difference
of $S_tD$ at rank $j\ge\ell+t+1$ is
$D_{j-1}-D_{j-t-1}\le0$. The extra interval $z^2[t-1]$ has
finished by that rank. Thus $P$ has first mode at most $\ell+t$,
and the same splice gives
\[
               \operatorname{firstmode}(\widehat D)
                         \le\min(m+1,\ell+t).                 \tag{3}
\]
No strictness of the rising prefix, and no identification of first
and last modes, is used here.

Consequently $n-1\le m$ and
\[
                       n-t\le\ell<\ell+1\le n-1.
\]
The exact source recurrence is
$\widehat E=S_tE+z^{t+1}$, whence
\[
 \Delta\widehat E_n
   =\sum_{i=n-t}^{n-1}\Delta E_i
                      +\mathbf1_{n=t+1}-\mathbf1_{n=t+2}.      \tag{4}
\]
Every term in this sum is nonnegative: for $3\le i\le m$ it is
$2D_{i-1}-D_i\ge0$ by actual single-arm doubling, and the
zero-extended ranks below three contribute zero. The contained rank
$i=\ell+1$ contributes
\[
               2D_\ell-D_{\ell+1}\ge D_\ell=M.
\]
Only one negative unit occurs in (4), proving (1).

At the old first missing rank the exact full-prefix identity is
\[
                 D_\kappa=2^{\kappa-2}-\delta.
\]
For $\kappa\ge8$ this is positive. The actual multiplicity bound
$\delta\le\kappa$ therefore gives
$M-1\ge2^{\kappa-2}-\kappa-1$.
For every integer $\kappa\ge8$,
\[
        2^{\kappa-2}-\kappa-1
                   \ge\frac{\kappa(\kappa+3)}2
                   =\binom{\kappa+2}{2}-1.                    \tag{5}
\]
Indeed twice the desired inequality is
$2\cdot2^{\kappa-2}\ge\kappa^2+5\kappa+2$.
It holds at eight. Doubling the polynomial right side dominates
its successor, since the difference is $\kappa^2+3\kappa-4\ge0$.
Induction proves (5). Now
$\widehat\delta\le\kappa+1$ proves (2).
\end{proof}

The pre-old-peak range and the peak-crossing cases with old
$\kappa\le7$ are now also closed by the complete theorem cited above.
A source-curvature bound alone does not settle later
source overlap: the growing convolution of the two sources still
requires its own payment. This result is not registered as a complete
both-rising S theorem.
